\section{Triadic Euler extension and preservation of forms}
\label{eul:capitulo}

The quadratic representation of TRIT acquires operator content by
identifying the generators selected by APP and by the TPK
calendar. This chapter reconstructs those antecedents before expressing their
exponential realizations. The order distinguishes the discrete operation,
its matrix representation and the analytic function that allows its
evolution in a real parameter to be described.

Three constructions are involved: the multiplicative cycle of order three
in APP, the elementary transport preserving a prolongation
coordinate and the incidence between modes descending from the tritic
calendar. Their generators satisfy the elliptic, parabolic and hyperbolic
relations, respectively. The alternating form of the plane then allows
a common conservation to be established, with the metric
and orientation differences of each realization explicitly preserved.

\subsection{APP multiplicative cycle and elliptic generator}
\label{eul:coxeter}

In \(\mathbb Z/9\mathbb Z\), multiplication by four is invertible
and has order three:
\[
 4^2=7\pmod9,\qquad 4^3=1\pmod9.
\]
On the units appear the cycles \((1,4,7)\) and \((2,8,5)\).
The classes \(0,3,6\) remain fixed. The selection of the multiplier and its
orbits therefore belongs to the nonadic arithmetic of APP.

\subsubsection{Contragredient action on the two evaluations}

Let \(\bar S(x,y)=x+y\) and \(\bar P(x,y)=xy\), with values in
\(\mathbb Z/9\mathbb Z\). For \(a=4\),
\begin{equation}
 \bar S(ax,ay)=a\bar S(x,y),\qquad
 \bar P(ax,ay)=a^2\bar P(x,y)=a^{-1}\bar P(x,y).
\label{eul:eq:contragrediente}
\end{equation}
The first identity follows from distributivity; the second, from
\(a^3=1\). Sum and product traverse opposite orientations of the
same residual cycle. In integer evaluations, the nonadic
quotients remain associated with the lift of these classes.

\subsubsection{Augmentation module and cycle matrix}

Fixing one of the cycles \((r,4r,7r)\), consider the free module
generated by \(e_r,e_{4r},e_{7r}\). Its augmentation submodule is
\[
 L=\{x_0e_r+x_1e_{4r}+x_2e_{7r}:x_0+x_1+x_2=0\}.
\]
A basis is \(b_1=e_r-e_{4r}\), \(b_2=e_{4r}-e_{7r}\). The cycle
sends \(b_1\) to \(b_2\) and \(b_2\) to \(-b_1-b_2\). Its matrix,
acting on column vectors, is
\[
 C=\begin{pmatrix}0&-1\\1&-1\end{pmatrix}.
\]
The scalar product inherited from the three-coordinate module has Gram
matrix
\[
 G_{\mathrm e}=\begin{pmatrix}2&-1\\-1&2\end{pmatrix}.
\]
One thus obtains the lattice \(A_2\), with its form determined by the
differences between positions of the cycle.

\begin{proposicion}[Elliptic normalization of the APP cycle]
\label{eul:prop:eliptico}
The cycle satisfies
\[
 C^2+C+I=0,\qquad C^\mathsf TG_{\mathrm e}C=G_{\mathrm e}.
\]
In \(L\otimes_{\mathbb Z}\mathbb R\), the operator
\[
 J_{+1}=\frac{2C+I}{\sqrt3}
\]
has trace zero and satisfies \(J_{+1}^2=-I\).
\end{proposicion}
\begin{proof}
Multiplication gives
\(C^2=\bigl(\begin{smallmatrix}-1&1\\-1&0\end{smallmatrix}\bigr)\),
whence \(C^2+C+I=0\). The cyclic permutation preserves the
scalar product of the three coordinates and its augmentation submodule;
in the chosen basis this is equivalent to the Gram identity. Finally,
\((2C+I)^2=4C^2+4C+I=-3I\), and
\(\operatorname{tr}(2C+I)=0\).
\end{proof}

Real normalization provides an operator root of \(-I\) from
the cycle already constructed. The identity can also be written
\[
 C=-\frac12I+\frac{\sqrt3}{2}J_{+1}.
\]
Its angular interpretation is obtained after fixing the analytic
realization of the elliptic type.

\subsection{Threshold transport and parabolic generator}
\label{eul:parabolico}

Consider the integer coordinates \((x,c)\), where \(c\) is preserved
as a recorded increment. The elementary transport
\[
 (x,c)\longmapsto(x+c,c)
\]
is represented by
\[
 U_0=I+N,
 \qquad N=\begin{pmatrix}0&1\\0&0\end{pmatrix}.
\]
This is the linear model of threshold transport used in the
quadratic realization. Balanced division and its normalization,
developed in chapter~\ref{trit:capitulo}, preserve their own
rules when applied to an integer evaluation.

\begin{proposicion}[Iteration and reversibility of transport]
\label{eul:prop:parabolico}
One has \(N^2=0\), \(N\ne0\) and, for every \(m\in\mathbb Z\),
\[
 U_0^m=I+mN,
 \qquad U_0^m(x,c)=(x+mc,c).
\]
Transport preserves the determinant and admits inverse \(I-N\).
\end{proposicion}
\begin{proof}
The image of \(N\) is contained in the first coordinate, on which
\(N\) vanishes; hence \(N^2=0\). The identity
\((I+aN)(I+bN)=I+(a+b)N\) proves the positive powers, the inverse and
the negative powers. The determinant of \(I+mN\) is one.
\end{proof}

The parabolic index \(\tau=0\) corresponds to a square-zero
generator. The operator \(N\) transports a linear datum and can be nonzero
on a particular state. The zero signature, the zero operator and a
nilpotent generator are three distinct objects.

\subsection{Mode incidence derived from the TPK calendar}
\label{eul:incidencia-modos}

The construction in section~\ref{sec:cal-incidencia} provides
the matrix~\eqref{eq:cal-incidencia}. Its primitive action is made explicit here
to link the tritic count with the hyperbolic quadratic representation
developed below.

The tritic phase selector takes, over the positions \(1,\ldots,9\),
the values \((+1,-1,0)\) repeated three times. On the operative mode
the rule used is
\[
 +1:\ m\mapsto\Sigma,\qquad
 -1:\ m\mapsto\Pi,\qquad
 0:\ m\mapsto m.
\]
After each update, the primitive block produces the sequence
of modes \((\Sigma,\Pi,\Pi)\). Spatial transport and register
updating remain in the complete TPK transition; here one reads
its modal incidence.

\begin{teorema}[Primitive incidence of the calendar]
\label{eul:teo:fav}
Counting consecutive pairs with cyclic closure in
\((\Sigma,\Pi,\Pi)\), the matrix obtained in the order \((\Sigma,\Pi)\)
is
\[
 F_{\rm av}=\begin{pmatrix}0&1\\1&1\end{pmatrix}.
\]
In calendars of length \(9,27,108\), the counts are
\(3F_{\rm av},9F_{\rm av},36F_{\rm av}\), respectively.
\end{teorema}
\begin{proof}
The three pairs are \(\Sigma\to\Pi\), \(\Pi\to\Pi\) and
\(\Pi\to\Sigma\). Each appears once. The pair
\(\Sigma\to\Sigma\) appears zero times. This determines the four
entries. Repeating the calendar multiplies all counts
by the number of primitive blocks, which is \(3,9,36\) in the indicated
cases.
\end{proof}

The sheet reader associates \(\Sigma\) with the areal reading and \(\Pi\) with
the volumetric reading. This assignment transports the resulting incidence to
the two reading types. The matrix results from the calendar; its spectral
analysis is performed once its entries have been constructed.

Counts of repeated calendars and matrix powers have
different functions. \(36F_{\rm av}\) counts the pairs traversed
during a calendar of \(108\) events. By contrast,
\(F_{\rm av}^{108}\) counts admissible chains of length \(108\)
in the incidence graph. The first object follows a chronology;
the second belongs to its adjacency hull.

\subsubsection{Hyperbolic realization of incidence}

Two incidence compositions produce
\[
 A=F_{\rm av}^2=\begin{pmatrix}1&1\\1&2\end{pmatrix}.
\]
Its coefficients count two-edge paths. One has
\(\det A=1\), \(\operatorname{tr}A=3\) and, by direct multiplication,
\(A^2-3A+I=0\).

\begin{proposicion}[Hyperbolic generator and invariant form]
\label{eul:prop:hiperbolico}
The operator
\[
 J_{-1}=\frac{2A-3I}{\sqrt5}
\]
has trace zero and square \(I\). The symmetric form
\[
 G_{\mathrm h}=\begin{pmatrix}2&1\\1&-2\end{pmatrix}
\]
has signature \((1,1)\) and satisfies
\(A^\mathsf TG_{\mathrm h}A=G_{\mathrm h}\).
\end{proposicion}
\begin{proof}
From \(A^2-3A+I=0\) one obtains
\((2A-3I)^2=4A^2-12A+9I=5I\). The trace of the numerator is zero.
Moreover,
\[
 G_{\mathrm h}A=\begin{pmatrix}3&4\\-1&-3\end{pmatrix},
 \qquad
 A^\mathsf TG_{\mathrm h}A
 =\begin{pmatrix}2&1\\1&-2\end{pmatrix}.
\]
The determinant of \(G_{\mathrm h}\) is \(-5\), so its two
eigenvalues have opposite signs.
\end{proof}

The realization is Lorentzian in dimension \(1+1\): a real plane,
a form of signature \((1,1)\) and a transport preserving it
are specified. The physical realization of a spacetime subsequently uses
its corresponding fields, dimensions and transduction maps.

\subsection{Exponential identity of the three regimes}
\label{eul:identidad}

The three concrete operators act in spaces of distinct provenance:
the augmentation module of an APP cycle, the threshold-transport plane
and the modal-incidence plane. Their normalizations satisfy
\[
 J_\tau^2=-\tau I,\qquad \operatorname{tr}J_\tau=0,
 \qquad \tau\in\{+1,0,-1\}.
\]
This relation determines their powers and allows the evolution
functions to be constructed by convergent series.

\begin{definicion}[Functions of the quadratic regime]
For \(t\in\mathbb R\), one defines
\[
 c_\tau(t)=\sum_{k=0}^{\infty}
       \frac{(-\tau)^kt^{2k}}{(2k)!},\qquad
 s_\tau(t)=\sum_{k=0}^{\infty}
       \frac{(-\tau)^kt^{2k+1}}{(2k+1)!}.
\]
These functions analytically realize the quadratic type already obtained.
\end{definicion}

\begin{teorema}[Triadic extension of Euler’s identity]
\label{eul:teo:euler}
For the preceding generators,
\begin{equation}
 E_\tau(t):=\exp(tJ_\tau)
 =c_\tau(t)I+s_\tau(t)J_\tau.
\label{eul:eq:euler}
\end{equation}
The specializations of \((c_\tau,s_\tau)\) are
\((\cos,\sin)\), \((1,t)\) and \((\cosh,\sinh)\), in the order
\(\tau=+1,0,-1\). Moreover,
\[
 c_\tau(t)^2+\tau s_\tau(t)^2=1,
 \qquad E_\tau(t)^{-1}=E_\tau(-t).
\]
\end{teorema}
\begin{proof}
The exponential series converges absolutely in matrix norm.
The quadratic relation gives
\(J_\tau^{2k}=(-\tau)^kI\) and
\(J_\tau^{2k+1}=(-\tau)^kJ_\tau\). Separating the even and odd
terms produces \eqref{eul:eq:euler}. The three specializations are
recognized in their series. Since the matrices \(tJ_\tau\) and \(-tJ_\tau\)
commute, their exponentials have product \(I\). The parity of the
series gives
\[
 E_\tau(t)E_\tau(-t)
 =\bigl(c_\tau(t)^2+\tau s_\tau(t)^2\bigr)I,
\]
which proves the scalar identity and the inverse.
\end{proof}

The numbers appearing when analytic functions are evaluated have here
the status of recognition of previously constructed operators.
Coinductive regional generation preserves its emitter, its lifts
and its cylinders. In particular, these series describe the exponential
realization; selection of the regional prefixes is performed through
the TPK.

\subsubsection{Composition law and quadratic invariants}

Commutation of multiples of the same generator gives
\(E_\tau(t)E_\tau(s)=E_\tau(t+s)\). Comparing coefficients of
\(I,J_\tau\),
\begin{align}
 c_\tau(t+s)&=c_\tau(t)c_\tau(s)-\tau s_\tau(t)s_\tau(s),\\
 s_\tau(t+s)&=s_\tau(t)c_\tau(s)+c_\tau(t)s_\tau(s).
\end{align}
Quadratic conjugation changes \(t\) to \(-t\) within the fixed
regime. Oriented inversion of a tritic sequence, for its part,
can change the regime indices and must be transported through the
corresponding TPK involution.

In the elliptic type, the preserved identity is \(a^2+b^2=1\).
In the hyperbolic type, \(a^2-b^2=1\), and the coordinates \(u=a+b\),
\(v=a-b\) satisfy \(uv=1\). Hyperbolic evolution has
\((u,v)=(e^t,e^{-t})\): expansion of one coordinate is accompanied
by contraction of the other. In the parabolic type,
\(E_0(t)=I+tN\), with linear transport and determinant one.

\subsubsection{Evaluations on the discrete operators}

The representation of the elliptic cycle gives
\[
 C=E_{+1}(2\pi/3),\qquad
 C^n=E_{+1}(2\pi n/3),\quad n\in\mathbb Z.
\]
In particular,
\[
 E_{+1}(\pi)+I=0,\qquad E_{+1}(2\pi)=I,
 \qquad E_0(1)=I+N.
\]
The angular coordinate \(\pi\) expresses the half-turn analytically;
its identification with the HMT regional output belongs to the coinductive
reader developed in its causal position.

The incidence \(F_{\rm av}\) has characteristic polynomial
\(x^2-x-1\). Its positive root \(\lambda\) satisfies
\(\lambda^2=\lambda+1\); recognizing it as the golden ratio, one writes
\(\lambda=\varphi\). With \(\eta=2\log\lambda\),
\[
 \cosh\eta=\frac32,\qquad \sinh\eta=\frac{\sqrt5}{2},
 \qquad E_{-1}(\eta)=A=F_{\rm av}^2.
\]
The check uses \(\lambda^{-1}=\lambda-1\), whence
\(\lambda^2+\lambda^{-2}=3\) and
\(\lambda^2-\lambda^{-2}=\sqrt5\). The matrix was obtained earlier by
counting transitions; spectral recognition characterizes its
scale. Equality with the generated regional coordinate is preserved
with the reader and the proof of that region.

Likewise, the scalar identity \(\exp(I)=eI\) characterizes the unit
 evaluation of the exponential. Its function in the analytic representation
is distinct from the coinductive construction of the corresponding region.

\subsection{Polynomial resolution of the regimes}
\label{eul:proyectores}

In the direct sum of the three real spaces, define
\[
 \mathbb J=J_{+1}\oplus N\oplus J_{-1},\qquad Q=\mathbb J^2.
\]
On those summands, \(Q\) acts as \(-I,0,I\), respectively.
Therefore \(Q^3=Q\) and \(\mathbb J^6=\mathbb J^2\).

\begin{proposicion}[Regime idempotents]
\label{eul:prop:proyectores}
The operators
\[
 p_{\mathrm e}=\frac{Q^2-Q}{2},\qquad
 p_{\mathrm p}=I-Q^2,\qquad
 p_{\mathrm h}=\frac{Q^2+Q}{2}
\]
are orthogonal idempotents in the algebraic sense and sum to \(I\).
Their images are the three summands of the representation.
\end{proposicion}
\begin{proof}
The three polynomials take the values \((1,0,0)\), \((0,1,0)\) and
\((0,0,1)\) on \(-1,0,+1\). Their squares minus themselves and
their cross products vanish at those three points, so they are
multiples of \(X^3-X\). Substituting \(Q\) demonstrates idempotence and
vanishing of the cross products. The sum of the three polynomials
is one. Their action on each block identifies the images.
\end{proof}

The complete exponential is expressed by
\begin{align}
 \exp(t\mathbb J)={}&p_{\mathrm e}(\cos t\,I+\sin t\,\mathbb J)
 +p_{\mathrm p}(I+t\mathbb J)\notag\\
 &+p_{\mathrm h}(\cosh t\,I+\sinh t\,\mathbb J).
\end{align}
This resolution distinguishes the regimes exactly. Transitions
between them belong to the TPK maps on states with orientation,
mode and memory; the direct sum keeps their domains separate.

\subsection{Nonadic compression and complementary register}
\label{eul:compresion}

Conservation between a reading and its complementary register follows
from the nine-phase structure. Let \(V\) be a space with a scalar
product and \(C:V\to V\) a linear transport. In \(V^9\), define
\[
 S_C(v_0,\ldots,v_8)=(Cv_8,v_0,\ldots,v_7),
 \qquad Jv=\frac13(v,\ldots,v).
\]
The inclusion \(J\) is isometric because it brings together nine copies with factor
\(1/3\). Its adjoint sums the components and divides by three. The
compression onto that inclusion is
\[
 T=J^*S_CJ=\frac{8I+C}{9}.
\]
Eight components preserve \(v\) and one passes through the transport \(C\).
The orthogonal complement of the inclusion records
\[
 (I-JJ^*)S_CJv
 =\frac1{27}\bigl(8(C-I)v,-(C-I)v,\ldots,-(C-I)v\bigr).
\]
The square of its norm is
\(8\|(C-I)v\|^2/81\). The same vector content can be represented
isometrically by
\[
 D=\frac{\sqrt8}{9}(C-I),
\]
maintaining the nine-component register when their
individual labels are used.

\begin{proposicion}[Algebraic recovery from two readings]
\label{eul:prop:recuperacion}
For every \(v\in V\),
\begin{equation}
 v=Tv-\frac1{\sqrt8}Dv.
\label{eul:eq:recuperacion}
\end{equation}
If \(C\) is isometric, one additionally has
\[
 \|v\|^2=\|Tv\|^2+\|Dv\|^2.
\]
\end{proposicion}
\begin{proof}
The first identity follows from
\((8I+C-(C-I))/9=I\). For the second, expanding the products
gives
\[
 T^*T+D^*D=\frac{72I+9C^*C}{81}.
\]
The hypothesis \(C^*C=I\) gives the identity. The mixed terms cancel with coefficients \(+8\) and \(-8\).
\end{proof}

Recovery uses the vector records, including their correlations. The intensities \(\|Tv\|^2\) and \(\|Dv\|^2\) provide the scalar balance, whereas the inverse formula uses the complete \(Tv\) and \(Dv\).

\subsection{Conservation of oriented area}
\label{eul:area}

In any oriented real plane, let
\[
 \Omega=\begin{pmatrix}0&1\\-1&0\end{pmatrix}.
\]
For a matrix \(L=\bigl(\begin{smallmatrix}a&b\\c&d\end{smallmatrix}\bigr)\), direct multiplication gives
\begin{equation}
 L^\mathsf T\Omega L=(ad-bc)\Omega=(\det L)\Omega.
\label{eul:eq:determinante-area}
\end{equation}
Since the tritic generators have zero trace, \(\det E_\tau(t)=1\). Their realizations therefore preserve the alternating form, although their metric forms differ.

\begin{teorema}[Nonadic conservation common to the three regimes]
\label{eul:teo:area} For \(C=E_\tau(t)\), let \(T=(8I+C)/9\) and \(D=\sqrt8(C-I)/9\). Then
\begin{equation}
 \Omega=T^\mathsf T\Omega T+D^\mathsf T\Omega D.
\label{eul:eq:balance-area}
\end{equation}
\end{teorema}
\begin{proof}
Expansion gives
\[
 T^\mathsf T\Omega T+D^\mathsf T\Omega D
 =\frac{72\Omega+9C^\mathsf T\Omega C}{81}.
\]
The terms \(C^\mathsf T\Omega\) and \(\Omega C\) cancel. By \eqref{eul:eq:determinante-area}, \(C^\mathsf T\Omega C=\Omega\), and the result is \(\Omega\).
\end{proof}

With \(c=c_\tau(t)\), the determinant contributions are
\begin{equation}
 \det T=\frac{65+16c}{81},\qquad
 \det D=\frac{16(1-c)}{81},\qquad
 \det T+\det D=1.
\label{eul:eq:determinantes}
\end{equation}
Indeed, for \(\det C=1\), \(\det(aI+C)=a^2+a\operatorname{tr}C+1\), and \(\operatorname{tr}C=2c\).

\begin{tablacompleta}
\caption{Contributions to the oriented area of three generated transports.}
\begin{tabular}{lccc}
\toprule
Transporte&\(\operatorname{tr}C\)&\(\det T\)&\(\det D\)\\
\midrule APP cycle \(C\)&$-1$&$19/27$&$8/27$\\
Threshold \(I+N\)&$2$&$1$&$0$\\
Incidence \(F_{\rm av}^2\)&$3$&$89/81$&$-8/81$\\
\bottomrule
\end{tabular}
\notatabla{Each pair sums to one. The sign records orientation; the table concerns alternating area, not probabilities.}
\end{tablacompleta}

In parabolic transport, \(D\) may be nonzero and have rank one: its area is zero and its record contains linear information. In hyperbolic transport, the complement has opposite orientation. In elliptic transport, the positive form intrinsic to the cycle also permits the norm balance. The three realizations share the alternating identity with their types preserved.

\subsection{Variable forms and restitution between successive planes}
\label{eul:formas-variables}

Consider a sequence of planes \(V_n\) with positive-definite symmetric matrices \(G_n\). Let \(B_n,C_n:V_n\to V_{n+1}\) be two transports satisfying
\[
 B_n^\mathsf TG_{n+1}B_n=G_n,
 \qquad C_n^\mathsf TG_{n+1}C_n=G_n.
\]
The transport \(B_n\) identifies the coordinates between the two planes; \(C_n\) includes the additional evolution. Define
\[
 T_n=\frac{8B_n+C_n}{9},\qquad
 D_n=\frac{\sqrt8}{9}(C_n-B_n).
\]

\begin{teorema}[Conservation and recovery with a variable form]
\label{eul:teo:forma-variable}
The following hold
\begin{align}
 G_n&=T_n^\mathsf TG_{n+1}T_n+D_n^\mathsf TG_{n+1}D_n,
 \label{eul:eq:balance-variable}\\
 v&=B_n^{-1}\left(T_nv-\frac1{\sqrt8}D_nv\right).
 \label{eul:eq:recuperacion-variable}
\end{align}
\end{teorema}
\begin{proof}
In the expansion of the first expression, the cross terms cancel. This leaves \(72B_n^\mathsf TG_{n+1}B_n/81\) and \(9C_n^\mathsf TG_{n+1}C_n/81\), whose sum is \(G_n\). The identity \(T_n-D_n/\sqrt8=B_n\) proves the second formula. Positivity of the forms makes \(B_n\) injective; between planes of equal dimension, it is invertible.
\end{proof}

For \(I_n(v)=\tfrac12v^\mathsf TG_nv\), \(v_{n+1}=T_nv_n\) and \(m_n=D_nv_n\), one obtains
\[
 I_n(v_n)=I_{n+1}(v_{n+1})+I_{n+1}(m_n).
\]
Summing over \(N\) steps gives
\begin{equation}
 I_0(v_0)=I_N(v_N)+\sum_{n=0}^{N-1}I_{n+1}(m_n).
\label{eul:eq:telescopica}
\end{equation}
Each record is evaluated in the form of the plane where it was produced. The inverse transports of \eqref{eul:eq:recuperacion-variable} restore all states in reverse order from the terminal state and the complete records.

A concrete realization uses invertible matrices \(M_n\), \(G_n=M_n^{-\mathsf T}M_n^{-1}\) and rotations \(R_n\):
\[
 B_n=M_{n+1}M_n^{-1},\qquad
 C_n=M_{n+1}R_nM_n^{-1}.
\]
The two form identities are verified by substituting these expressions and using \(R_n^\mathsf TR_n=I\). The coordinates, scales and parameters of \(M_n,R_n\) are received from the state readers when this realization is applied to Dimensional Mechanics.

\subsubsection{Exact criterion for transfer to the record}

If \(R_n\) is a rotation through angle \(\theta_n\), in normalized coordinates \(y_n=M_n^{-1}v_n\) one has
\[
 y_{n+1}=\frac{8I+R_n}{9}y_n,
 \qquad
 \left(\frac{8I+R_n}{9}\right)^\mathsf T
 \left(\frac{8I+R_n}{9}\right)=r_nI,
\]
\[
 r_n=1-\frac{32}{81}\sin^2(\theta_n/2).
\]
By induction, \(I_N(v_N)=I_0(v_0)\prod_{n<N}r_n\).

\begin{proposicion}[Criterion for vanishing of the terminal state]
If \(I_0(v_0)>0\), the terminal state satisfies \(I_N(v_N)\to0\) exactly when
\[
 \sum_{n\ge0}\sin^2(\theta_n/2)=+\infty.
\]
\end{proposicion}
\begin{proof}
For \(c=32/81\) and \(u\in[0,1]\), integrating \(1/(1-s)\) between zero and \(cu\) gives
\[
 cu\le-\log(1-cu)\le\frac{cu}{1-c}.
\]
The sum of the logarithms diverges exactly when the sum of the \(u_n=\sin^2(\theta_n/2)\) diverges. The product of the \(r_n\) tends to zero exactly in that case.
\end{proof}

The conservation law \eqref{eul:eq:telescopica} holds at every finite depth. The preceding criterion determines when the entire initial form is transferred to the record in this realization. The angular sequence is the one obtained by the corresponding trajectory reader.

\subsection{Integrated action and compatibility with refinement}
\label{eul:accion-refinamiento}

In a canonical plane with coordinates \((Q,P)\), let \(\gamma\) be a piecewise \(C^1\) closed curve. For any real matrix \(L\) of order two,
\begin{equation}
 \oint_{L\gamma}P\,\mathrm dQ
 =\det(L)\oint_\gamma P\,\mathrm dQ.
\label{eul:eq:accion-determinante}
\end{equation}
To prove this, write \(Q'=aQ+bP\), \(P'=cQ+dP\). The closed integrals of \(Q\,\mathrm dQ\) and \(P\,\mathrm dP\) are zero. Integrating \(\mathrm d(PQ)\) gives \(\oint Q\,\mathrm dP=-\oint P\,\mathrm dQ\). The coefficient \(ad-bc\) remains, proving the formula.

Applied to \eqref{eul:eq:determinantes}, this yields
\[
 \oint_\gamma P\,\mathrm dQ
 =\oint_{T\gamma}P\,\mathrm dQ
  +\oint_{D\gamma}P\,\mathrm dQ.
\]
The orientation of each curve preserves the sign of its contribution.

Refinement additionally requires compatibility between operators at successive levels. In this section, let \(\iota_n:V_n\to V_{n+1}\) be linear injections and \(A_n:V_n\to V_n\) endomorphisms satisfying \(A_{n+1}\iota_n=\iota_nA_n\). At each level we define \(T_n=(8I_{V_n}+A_n)/9\) and \(D_n=\sqrt8(A_n-I_{V_n})/9\). These expressions imply
\[
 T_{n+1}\iota_n=\iota_nT_n,
 \qquad D_{n+1}\iota_n=\iota_nD_n.
\]
Restitution \eqref{eul:eq:recuperacion} then commutes with refinement. This identity preserves the operation and the record simultaneously; its verification in a concrete realization uses the map \(\iota_n\) produced by that prolongation.

\subsubsection{Preservation of the reader's scalar factor}

If a realization receives a scale factor \(e^{-x}\), the transport is \(C=e^{-x}E_\tau(t)\). Its determinant is \(e^{-2x}\), and the same expansion of the nonadic law gives
\[
 T^\mathsf T\Omega T+D^\mathsf T\Omega D
 =\frac{8+e^{-2x}}9\Omega.
\]
For \(x\ge0\), the scale complement
\[
 D_{\rm esc}=\frac{\sqrt{1-e^{-2x}}}{3}I
\]
restores the complete balance:
\[
 \Omega=T^\mathsf T\Omega T+D^\mathsf T\Omega D
       +D_{\rm esc}^\mathsf T\Omega D_{\rm esc}.
\]
The coefficient is obtained from the determinant defect of the scale. Each application preserves its reader's factor and specifies whether this additional record is needed.

The triadic extension of Euler is thus articulated through generated operations, invariant forms and recovery operators. The APP cycle, nilpotent transport and modal incidence preserve their provenance; the analytic realization describes their evolutions; the nonadic division of the reading preserves the complement required for restitution. This sequence expresses the link between regime, transformation and memory without replacing the genealogy with an isolated exponential formula.
